% ============================================================
%  TD Chapitre 5 : Mesures des angles
% ============================================================
\chapter{Mesures des angles}
\chaptermark{Mesures des angles}

\begin{exercice}[conversions degrés-radians]
Convertir en radians les trois angles $45\dgr{}$, $210\dgr{}$ et $300\dgr{}$, en simplifiant chaque fraction. Convertir ensuite en degrés les trois angles $\dfrac{5\pi}{4}$ rad, $\dfrac{11\pi}{6}$ rad et $\dfrac{\pi}{9}$ rad. Contrôler chaque résultat avec la correspondance $180\dgr{} = \pi$ rad.
\end{exercice}

\begin{exercice}[l'angle entre les aiguilles d'une horloge]
Chaque heure vaut $30\dgr{}$ et chaque minute $6\dgr{}$ sur le cadran. Déterminer l'angle entre la petite aiguille et la grande aiguille à $8$ h $20$, en :
\begin{enumerate}
\item calculant la position de la petite aiguille depuis midi (attention : à $8$ h $20$, elle a dépassé le $8$ d'un tiers d'heure) ;
\item calculant la position de la grande aiguille ;
\item en déduisant l'angle entre les deux aiguilles.
\end{enumerate}
\end{exercice}

\begin{exercice}[Pythagore et trigonométrie]
\begin{enumerate}
\item une échelle de $6{,}5$ m est appuyée contre un mur, son pied étant à $2{,}5$ m du mur. À quelle hauteur le sommet de l'échelle touche-t-il le mur ?
\item un technicien se place à $25$ m du pied d'une tour ; l'angle entre le sol horizontal et la ligne de visée du sommet vaut $30\dgr{}$. Calculer la hauteur de la tour, avec $\tan 30\dgr{} \approx 0{,}577$.
\end{enumerate}
\end{exercice}

\begin{exercice}[opérations sur les vecteurs du plan]
Soient $\vec{u} = \begin{pmatrix} 4 \\ -3 \end{pmatrix}$ et $\vec{v} = \begin{pmatrix} 1 \\ 5 \end{pmatrix}$.
\begin{enumerate}
\item calculer $\vec{u} + \vec{v}$, puis les normes $\|\vec{u}\|$ et $\|\vec{v}\|$ (la seconde sera donnée sous la forme $a\sqrt{b}$ simplifiée puis arrondie à $0{,}01$ près) ;
\item calculer $2\vec{u} - \vec{v}$ et sa norme, arrondie à $0{,}01$ près ;
\item calculer la distance entre les points $A(-2, 4)$ et $B(3, -8)$.
\end{enumerate}
\end{exercice}
